\begin{thebibliography}{99}
\bibitem{rfc6550} T. Winter, P. Thubert, A. Brandt, J. Hui, R. Kelsey, P. Levis, K. Pister, R. Struik, J. P. Vasseur, and R. Alexander, ``RPL: IPv6 Routing Protocol for Low-Power and Lossy Networks,'' Internet Engineering Task Force, RFC 6550, Mar. 2012, doi: \href{https://doi.org/10.17487/RFC6550}{10.17487/RFC6550}.
\bibitem{albinali2024} H. Albinali and F. Azzedin, ``Toward RPL attacks and mitigation taxonomy: Systematic literature review approach,'' \emph{IEEE Trans. Netw. Serv. Manag.}, vol. 21, no. 5, pp. 5215--5238, 2024, doi: \href{https://doi.org/10.1109/TNSM.2024.3386468}{10.1109/TNSM.2024.3386468}.
\bibitem{raoof2019} A. Raoof, A. Matrawy, and C.-H. Lung, ``Routing attacks and mitigation methods for RPL-based Internet of Things,'' \emph{IEEE Commun. Surveys Tuts.}, vol. 21, no. 2, pp. 1582--1606, 2019, doi: \href{https://doi.org/10.1109/COMST.2018.2885894}{10.1109/COMST.2018.2885894}.
\bibitem{rfc6206} P. Levis, T. Clausen, J. Hui, O. Gnawali, and J. Ko, ``The Trickle algorithm,'' Internet Engineering Task Force, RFC 6206, Mar. 2011, doi: \href{https://doi.org/10.17487/RFC6206}{10.17487/RFC6206}.
\bibitem{rajasekar2022} V. R. Rajasekar and S. Rajkumar, ``A study on impact of DIS flooding attack on RPL-based 6LowPAN network,'' \emph{Microprocess. Microsyst.}, vol. 94, Art. no. 104675, 2022, doi: \href{https://doi.org/10.1016/j.micpro.2022.104675}{10.1016/j.micpro.2022.104675}.
\bibitem{vermaranga2020} A. Verma and V. Ranga, ``Mitigation of DIS flooding attacks in RPL-based 6LoWPAN networks,'' \emph{Trans. Emerg. Telecommun. Technol.}, vol. 31, no. 2, Art. no. e3802, 2020, doi: \href{https://doi.org/10.1002/ett.3802}{10.1002/ett.3802}.
\bibitem{medjek2021} F. Medjek, D. Tandjaoui, N. Djedjig, and I. Romdhani, ``Multicast DIS attack mitigation in RPL-based IoT-LLNs,'' \emph{J. Inf. Secur. Appl.}, vol. 61, Art. no. 102939, 2021, doi: \href{https://doi.org/10.1016/j.jisa.2021.102939}{10.1016/j.jisa.2021.102939}.
\bibitem{rfc6554} J. Hui, J. P. Vasseur, D. Culler, and V. Manral, ``An IPv6 routing header for source routes with the Routing Protocol for Low-Power and Lossy Networks (RPL),'' Internet Engineering Task Force, RFC 6554, Mar. 2012, doi: \href{https://doi.org/10.17487/RFC6554}{10.17487/RFC6554}.
\bibitem{pu2018} C. Pu, ``Mitigating DAO inconsistency attack in RPL-based low power and lossy networks,'' in \emph{Proc. IEEE 8th Annu. Comput. Commun. Workshop Conf. (CCWC)}, 2018, pp. 570--574, doi: \href{https://doi.org/10.1109/CCWC.2018.8301614}{10.1109/CCWC.2018.8301614}.
\bibitem{ghaleb2019} B. Ghaleb, A. Al-Dubai, E. Ekonomou, M. Qasem, I. Romdhani, and L. Mackenzie, ``Addressing the DAO insider attack in RPL's Internet of Things networks,'' \emph{IEEE Commun. Lett.}, vol. 23, no. 1, pp. 68--71, 2019, doi: \href{https://doi.org/10.1109/LCOMM.2018.2878151}{10.1109/LCOMM.2018.2878151}.
\bibitem{limsd2024} A. Verma, S. K. Verma, A. C. Pandey, J. Grover, and G. Sharma, ``Li-MSD: A lightweight mitigation solution for DAO insider attack in RPL-based IoT,'' \emph{Future Gener. Comput. Syst.}, vol. 159, pp. 327--339, 2024, doi: \href{https://doi.org/10.1016/j.future.2024.05.032}{10.1016/j.future.2024.05.032}.
\bibitem{prajisha2025} C. Prajisha and A. R. Vasudevan, ``Detection and mitigation of multicast DIS flooding attacks in RPL-based IoT networks,'' \emph{Ad Hoc Netw.}, vol. 179, Art. no. 104002, 2025, doi: \href{https://doi.org/10.1016/j.adhoc.2025.104002}{10.1016/j.adhoc.2025.104002}.
\bibitem{ganeriwal2008} S. Ganeriwal, L. K. Balzano, and M. B. Srivastava, ``Reputation-based framework for high integrity sensor networks,'' \emph{ACM Trans. Sensor Netw.}, vol. 4, no. 3, pp. 1--37, 2008, doi: \href{https://doi.org/10.1145/1362542.1362546}{10.1145/1362542.1362546}.
\bibitem{watkins1992} C. J. C. H. Watkins and P. Dayan, ``Q-learning,'' \emph{Mach. Learn.}, vol. 8, nos. 3--4, pp. 279--292, 1992, doi: \href{https://doi.org/10.1007/BF00992698}{10.1007/BF00992698}.
\bibitem{khan2023} M. A. Khan, R. N. B. Rais, and O. Khalid, ``Collaborative detection and prevention of Sybil attacks against RPL-based Internet of Things,'' \emph{Comput. Mater. Continua}, vol. 77, no. 1, pp. 827--843, 2023, doi: \href{https://doi.org/10.32604/cmc.2023.040756}{10.32604/cmc.2023.040756}.
\bibitem{sharma2022} D. K. Sharma, S. K. Dhurandher, S. Kumaram, K. Datta Gupta, and P. K. Sharma, ``Mitigation of black hole attacks in 6LoWPAN RPL-based wireless sensor network for cyber physical systems,'' \emph{Comput. Commun.}, vol. 189, pp. 182--192, 2022, doi: \href{https://doi.org/10.1016/j.comcom.2022.04.003}{10.1016/j.comcom.2022.04.003}.
\bibitem{ankam2023} S. Ankam and N. S. Reddy, ``A mechanism to detecting flooding attacks in quantum enabled cloud-based low-power and lossy networks,'' \emph{Theor. Comput. Sci.}, vol. 941, pp. 29--38, 2023, doi: \href{https://doi.org/10.1016/j.tcs.2022.08.018}{10.1016/j.tcs.2022.08.018}.
\bibitem{abhinaya2021} E. V. Abhinaya and B. Sudhakar, ``A secure routing protocol for low power and lossy networks based 6LoWPAN networks to mitigate DIS flooding attacks,'' \emph{J. Ambient Intell. Humaniz. Comput.}, published online Jan. 2021, doi: \href{https://doi.org/10.1007/s12652-020-02804-3}{10.1007/s12652-020-02804-3}.
\bibitem{nishanth2021} N. Nishanth and A. Mujeeb, ``Modeling and detection of flooding-based denial of service attacks in wireless ad hoc networks using uncertain reasoning,'' \emph{IEEE Trans. Cogn. Commun. Netw.}, vol. 7, no. 3, pp. 893--904, 2021, doi: \href{https://doi.org/10.1109/TCCN.2021.3055503}{10.1109/TCCN.2021.3055503}.
\bibitem{rouissat2023} M. Rouissat, M. Belkheir, I. S. Alsukayti, and A. Mokaddem, ``A lightweight mitigation approach against a new inundation attack in RPL-based IoT networks,'' \emph{Appl. Sci.}, vol. 13, no. 18, Art. no. 10366, 2023, doi: \href{https://doi.org/10.3390/app131810366}{10.3390/app131810366}.
\bibitem{imdrpl2023} P. Deepavathi and C. Mala, ``IMDRPL: Identifying and eliminating malicious devices using DIO and DAO ICMP control messages in RPL-based protocol,'' \emph{Peer-to-Peer Netw. Appl.}, vol. 16, no. 5, pp. 2380--2398, 2023, doi: \href{https://doi.org/10.1007/s12083-023-01539-0}{10.1007/s12083-023-01539-0}.
\bibitem{alsukayti2023} I. S. Alsukayti and M. Alreshoodi, ``RPL-based IoT networks under simple and complex routing security attacks: An experimental study,'' \emph{Appl. Sci.}, vol. 13, no. 8, Art. no. 4878, 2023, doi: \href{https://doi.org/10.3390/app13084878}{10.3390/app13084878}.
\bibitem{de2rarpl2024} P. Deepavathi and C. Mala, ``DE2RA-RPL: Detection and elimination of resource-related attacks in IoT RPL-based protocol,'' \emph{J. Supercomput.}, vol. 80, no. 15, pp. 22397--22427, 2024, doi: \href{https://doi.org/10.1007/s11227-024-06256-w}{10.1007/s11227-024-06256-w}.
\bibitem{abdulkareem2024} O. A. Abdulkareem, R. K. Kontham, and F. E. Mahmood, ``Collaborative intrusion detection system to identify joint attacks in routing protocol for low-power and lossy networks routing protocol on the Internet of Everything,'' \emph{Mesopotamian J. CyberSecurity}, vol. 4, no. 3, pp. 251--277, 2024, doi: \href{https://doi.org/10.58496/MJCS/2024/026}{10.58496/MJCS/2024/026}.
\bibitem{garciaribera2022} E. Garcia Ribera, B. Martinez Alvarez, C. Samuel, P. P. Ioulianou, and V. G. Vassilakis, ``An intrusion detection system for RPL-based IoT networks,'' \emph{Electronics}, vol. 11, no. 23, Art. no. 4041, 2022, doi: \href{https://doi.org/10.3390/electronics11234041}{10.3390/electronics11234041}.
\bibitem{flsecrpl2024} C. Kim, C. So-In, Y. Kongsorot, and P. Aimtongkham, ``FLSec-RPL: A fuzzy logic-based intrusion detection scheme for securing RPL-based IoT networks against DIO neighbor suppression attacks,'' \emph{Cybersecurity}, vol. 7, no. 1, Art. no. 27, 2024, doi: \href{https://doi.org/10.1186/s42400-024-00223-x}{10.1186/s42400-024-00223-x}.
\bibitem{dsatids2024} V. R. Rajasekar and S. Rajkumar, ``DSAT-IDS: Dissimilarity and adaptive threshold-based intrusion detection system to mitigate selective forwarding attack in the RPL-based 6LoWPAN,'' \emph{Cluster Comput.}, vol. 27, no. 9, pp. 12029--12068, 2024, doi: \href{https://doi.org/10.1007/s10586-024-04564-7}{10.1007/s10586-024-04564-7}.
\bibitem{azzaoui2024} H. Azzaoui, A. Z. E. Boukhamla, P. Perazzo, M. Alazab, and V. Ravi, ``A lightweight cooperative intrusion detection system for RPL-based IoT,'' \emph{Wireless Pers. Commun.}, vol. 134, no. 4, pp. 2235--2258, 2024, doi: \href{https://doi.org/10.1007/s11277-024-11009-2}{10.1007/s11277-024-11009-2}.
\bibitem{monteiro2026} B. Monteiro and J. Granjal, ``Trust-aware distributed and hybrid intrusion detection for rank attacks in RPL IoT environments,'' \emph{IoT}, vol. 7, no. 1, Art. no. 4, 2026, doi: \href{https://doi.org/10.3390/iot7010004}{10.3390/iot7010004}.
\bibitem{qsecrpl2023} G. Sharma, J. Grover, and A. Verma, ``QSec-RPL: Detection of version number attacks in RPL based mobile IoT using Q-learning,'' \emph{Ad Hoc Netw.}, vol. 142, Art. no. 103118, 2023, doi: \href{https://doi.org/10.1016/j.adhoc.2023.103118}{10.1016/j.adhoc.2023.103118}.
\bibitem{itrpl2024} D. Dey and N. Ghosh, ``iTRPL: An intelligent and trusted RPL protocol based on multi-agent reinforcement learning,'' \emph{Ad Hoc Netw.}, vol. 163, Art. no. 103586, 2024, doi: \href{https://doi.org/10.1016/j.adhoc.2024.103586}{10.1016/j.adhoc.2024.103586}.
\bibitem{haqdad2026} M. Haqdad, M. Fayaz, P. Khan, F. Ali, T. H. H. Aldhyani, A. K. Bashir, and D. Kwak, ``Intent-based networking with deep reinforcement learning for detecting decreased rank attacks in low-power and lossy IoT networks,'' \emph{IEEE Internet Things J.}, vol. 13, no. 9, pp. 17802--17814, 2026, doi: \href{https://doi.org/10.1109/JIOT.2025.3596097}{10.1109/JIOT.2025.3596097}.
\bibitem{ahmadi2026} K. Ahmadi and R. Javidan, ``An intelligent distributed mitigation approach for multiple routing attacks in RPL-based IoT networks,'' \emph{Wireless Netw.}, vol. 32, no. 5, pp. 2933--2957, 2026, doi: \href{https://doi.org/10.1007/s11276-026-04207-2}{10.1007/s11276-026-04207-2}.
\bibitem{panda2026} N. Panda and S. Muthuraman, ``A multiclass machine learning framework for detecting routing attacks in RPL-based IoT networks using a novel simulation-driven dataset,'' \emph{Future Internet}, vol. 18, no. 1, Art. no. 35, 2026, doi: \href{https://doi.org/10.3390/fi18010035}{10.3390/fi18010035}.
\bibitem{gyawali2021} S. Gyawali, Y. Qian, and R. Q. Hu, ``Deep reinforcement learning based dynamic reputation policy in 5G based vehicular communication networks,'' \emph{IEEE Trans. Veh. Technol.}, vol. 70, no. 6, pp. 6136--6146, 2021, doi: \href{https://doi.org/10.1109/TVT.2021.3079379}{10.1109/TVT.2021.3079379}.
\bibitem{jasti2024} P. C. Jasti, H. Griffith, and H. Rathore, ``A multi-level Dempster-Shafer and reinforcement learning-based reputation system for connected vehicle security,'' in \emph{Proc. IEEE 21st Consum. Commun. Netw. Conf. (CCNC)}, 2024, pp. 714--720, doi: \href{https://doi.org/10.1109/CCNC51664.2024.10454856}{10.1109/CCNC51664.2024.10454856}.
\bibitem{sentz2002} K. Sentz and S. Ferson, ``Combination of evidence in Dempster--Shafer theory,'' Sandia National Laboratories, Albuquerque, NM, USA, Tech. Rep. SAND2002-0835, 2002, doi: \href{https://doi.org/10.2172/800792}{10.2172/800792}.
\bibitem{sterbenz2013} J. P. G. Sterbenz, E. K. {\c{C}}etinkaya, M. A. Hameed, A. Jabbar, S. Qian, and J. P. Rohrer, ``Evaluation of network resilience, survivability, and disruption tolerance: Analysis, topology generation, simulation, and experimentation,'' \emph{Telecommun. Syst.}, vol. 52, no. 2, pp. 705--736, 2013, doi: \href{https://doi.org/10.1007/s11235-011-9573-6}{10.1007/s11235-011-9573-6}.
\bibitem{contiking2022} G. Oikonomou, S. Duquennoy, A. Elsts, J. Eriksson, Y. Tanaka, and N. Tsiftes, ``The Contiki-NG open source operating system for next generation IoT devices,'' \emph{SoftwareX}, vol. 18, Art. no. 101089, 2022, doi: \href{https://doi.org/10.1016/j.softx.2022.101089}{10.1016/j.softx.2022.101089}.
\bibitem{cooja2006} F. Osterlind, A. Dunkels, J. Eriksson, N. Finne, and T. Voigt, ``Cross-level sensor network simulation with COOJA,'' in \emph{Proc. 31st IEEE Conf. Local Comput. Netw. (LCN)}, 2006, pp. 641--648, doi: \href{https://doi.org/10.1109/LCN.2006.322172}{10.1109/LCN.2006.322172}.
\bibitem{rfc6719} O. Gnawali and P. Levis, ``The minimum rank with hysteresis objective function,'' Internet Engineering Task Force, RFC 6719, Sep. 2012, doi: \href{https://doi.org/10.17487/RFC6719}{10.17487/RFC6719}.
\bibitem{camp2002} T. Camp, J. Boleng, and V. Davies, ``A survey of mobility models for ad hoc network research,'' \emph{Wireless Commun. Mobile Comput.}, vol. 2, no. 5, pp. 483--502, 2002, doi: \href{https://doi.org/10.1002/wcm.72}{10.1002/wcm.72}.
\bibitem{wilcoxon1945} F. Wilcoxon, ``Individual comparisons by ranking methods,'' \emph{Biometrics Bull.}, vol. 1, no. 6, pp. 80--83, 1945, doi: \href{https://doi.org/10.2307/3001968}{10.2307/3001968}.
\bibitem{aickin1996} M. Aickin and H. Gensler, ``Adjusting for multiple testing when reporting research results: The Bonferroni vs Holm methods,'' \emph{Amer. J. Public Health}, vol. 86, no. 5, pp. 726--728, 1996, doi: \href{https://doi.org/10.2105/AJPH.86.5.726}{10.2105/AJPH.86.5.726}.
\bibitem{efron1994} B. Efron and R. J. Tibshirani, \emph{An Introduction to the Bootstrap}. New York, NY, USA: Chapman \& Hall/CRC, 1994, doi: \href{https://doi.org/10.1201/9780429246593}{10.1201/9780429246593}.
\end{thebibliography}
